% Unit 074 English translation of chapter6.tex lines 557--739.
% Source: Methods of Algebra, Volume 2, commit
% 9a5803ff77dd3257484cb177f851a73770a59dd3 (CC BY 4.0).
% stable-unit-id: o014.aljabr2.chapter6.low-degree-cohomology
% Coverage: Section 6.3 in full.

% segment-id: o014.li2.chapter6.low-degree-cohomology.h001
\section{Low-degree cohomology: crossed homomorphisms and group extensions}
\label{sec:grp-low-degree}
\hypertarget{o014.aljabr2.chapter6.low-degree-cohomology}{ }

% segment-id: o014.li2.chapter6.low-degree-cohomology.p001
This section aims to deepen our understanding of $\Hm^1$ and $\Hm^2$.
Fix a group $G$ and a commutative ring $\Bbbk$. Given a $G$-module $M$,
define the following submodules of $C^n(G,M)$ for the standard complex
$C(G,M)$ of Proposition~\ref{prop:groupcoh-cochain}:
\[ Z^n(G,M):=\Ker(d^n)\supset\Image(d^{n-1})=:B^n(G,M), \]
and set $B^0(G,M)=0$. We identify
\begin{equation*}
	\Hm^n(G,M)\quad\text{with}\quad Z^n(G,M)/B^n(G,M).
\end{equation*}

% segment-id: o014.li2.chapter6.low-degree-cohomology.p002
In view of their topological origins, the elements of $Z^n(G,M)$ are also
called $n$-\emph{cocycles}, while the elements of $B^n(G,M)$ are called
$n$-\emph{coboundaries}.
\index{cocycle}
\index{coboundary}
\index[sym1]{ZGM@$Z^n(G,M)$}
\index[sym1]{BGM@$B^n(G,M)$}

% segment-id: o014.li2.chapter6.low-degree-cohomology.p003
Recall that $C^n(G,M)$ consists of all maps $f:G^n\to M$. The low-degree
part of the standard complex is
\[\begin{tikzcd}[row sep=tiny, column sep=small]
	C^0(G,M) \arrow[r,"{d^0}"] & C^1(G,M) \arrow[r,"{d^1}"] & C^2(G,M) \\
	m \arrow[mapsto,r] & {[g\mapsto gm-m]} & \\
	& f \arrow[mapsto,r] &
	{[(g_1,g_2)\mapsto g_1f(g_2)-f(g_1g_2)+f(g_1)]},
\end{tikzcd}\]
In particular, $Z^0(G,M)=M^G$. We next consider the case $n=1$.

% segment-id: o014.li2.chapter6.low-degree-cohomology.d001
\begin{definition}[Crossed homomorphism]\label{def:crossed-homomorphism}
	\index{crossed homomorphism}
	Let $M$ be a $G$-module. The elements of $Z^1(G,M)$ are also called
	\emph{crossed homomorphisms} from $G$ to $M$; in other words, a crossed
	homomorphism is a map $f:G\to M$ satisfying
	\[ f(g_1g_2)=g_1f(g_2)+f(g_1),\qquad g_1,g_2\in G. \]
	For fixed $\Bbbk$, they form a $\Bbbk$-module under pointwise operations.
\end{definition}

% segment-id: o014.li2.chapter6.low-degree-cohomology.p004
Notice that the elements of $B^1(G,M)$ are the crossed homomorphisms of the
form $f_m:g\mapsto gm-m$, where $m\in M$.

% segment-id: o014.li2.chapter6.low-degree-cohomology.p005
In general, every crossed homomorphism $f$ satisfies $f(1_G)=0$ (take
$g_1=g_2=1_G$) and $g^{-1}f(g)=-f(g^{-1})$ (take $g_1=g^{-1}$ and $g_2=g$).

% segment-id: o014.li2.chapter6.low-degree-cohomology.e001
\begin{example}
	If $G$ acts trivially on $M$, then the crossed homomorphisms are precisely
	the homomorphisms from $G$ to the additive group $(M,+)$, and
	$B^1(G,M)=\{0\}$. Thus in this case
	\[ \Hom_{\cate{Grp}}(G,M)\simeq\Hm^1(G,M). \]
\end{example}

% segment-id: o014.li2.chapter6.low-degree-cohomology.p006
Another way to understand crossed homomorphisms is through semidirect products.
For a $G$-module $M$, use the $G$-action to form the semidirect product
$M\rtimes G$. Its elements have the form $(m,g)$ with $m\in M$ and $g\in G$,
and it has a projection homomorphism $\pi:M\rtimes G\to G$ sending $(m,g)$
to $g$.

% segment-id: o014.li2.chapter6.low-degree-cohomology.t001
\begin{proposition}\label{prop:crossed-semidirect-product}
	% Editorial correction: the Chinese source has Z^1(M,G); the arguments
	% and the Indonesian witness require Z^1(G,M).
	Let $M$ be a $G$-module. There is a canonical bijection
	\[ Z^1(G,M)\xrightarrow{1:1}
	\left\{\text{group homomorphisms }\sigma:G\to M\rtimes G
	\;\middle|\;\pi\sigma=\identity_G\right\}. \]
	A map $\sigma$ on the right-hand side is also called a section of $\pi$;
	the crossed homomorphism $f$ corresponds to the section
	$\sigma(g)=(f(g),g)$.
\end{proposition}
\begin{proof}
	A map $\sigma$ satisfying $\pi\sigma=\identity_G$ must have the form
	$\sigma(g)=(f(g),g)$ for some map $f:G\to M$. By the definition of the
	semidirect product, $\sigma$ is a group homomorphism if and only if the
	following equality holds in $M$:
	\[ f(g_1)+g_1f(g_2)=f(g_1g_2),\qquad g_1,g_2\in G. \]
	This is exactly the definition of a crossed homomorphism.
\end{proof}

% segment-id: o014.li2.chapter6.low-degree-cohomology.p007
Crossed homomorphisms are also related to the augmentation ideal
$\mathfrak{I}\subset\Bbbk[G]$ introduced in
Remark~\ref{rem:augmentation-kG}.

% segment-id: o014.li2.chapter6.low-degree-cohomology.t002
\begin{proposition}\label{prop:crossed-augmentation}
	For every $G$-module $M$, there is a canonical isomorphism
	$\Hom_G(\mathfrak{I},M)\rightiso Z^1(G,M)$ sending a homomorphism
	$\varphi$ to the crossed homomorphism $f(g)=\varphi(g-1_G)$.
\end{proposition}
\begin{proof}
	As a $\Bbbk$-module, $\mathfrak{I}$ has basis
	$\{g-1_G:g\in G,\;g\neq1_G\}$. Thus specifying a $\Bbbk$-module
	homomorphism $\varphi:\mathfrak{I}\to M$ is equivalent to specifying a
	map $f:G\to M$ satisfying $f(1_G)=0$, with
	$f(g)=\varphi(g-1_G)$. The condition that $\varphi$ be a $G$-module
	homomorphism is
	$\varphi(g_1(g_2-1_G))=g_1\varphi(g_2-1_G)$. But
	\[ \text{left-hand side}=\varphi(g_1g_2-1_G-g_1+1_G)
	=f(g_1g_2)-f(g_1),\qquad
	\text{right-hand side}=g_1f(g_2). \]
	This is exactly the crossed-homomorphism condition.
\end{proof}

% segment-id: o014.li2.chapter6.low-degree-cohomology.p008
For the next step, the case $n=2$, we first recall the notion of a group
extension. Take $\Bbbk=\Z$ and introduce some terminology.

% segment-id: o014.li2.chapter6.low-degree-cohomology.d002
\begin{definition}\label{def:group-ext}
	\index{group extension}
	Let $A$ be a group. An \emph{extension} of $G$ by $A$ is a short exact
	sequence of groups \cite[Definition 4.3.7]{Li1}
	\[ 0\to A\to E\xrightarrow{\pi}G\to1. \]

	Henceforth we abbreviate the data of the group extension to $E$. A group
	homomorphism $s:G\to E$ satisfying $\pi s=\identity_G$ is called a
	\emph{splitting} of the extension. In this case $s$ embeds $G$ as a
	subgroup of $E$, thereby identifying $E$ with the semidirect product
	$A\rtimes G$. An extension possessing a splitting is called
	\emph{split}. See \cite[\S 4.3]{Li1} for details.
\end{definition}

% segment-id: o014.li2.chapter6.low-degree-cohomology.p009
Broadly speaking, the theory of group extensions studies how to build a larger
group from a normal subgroup and a quotient group. To relate this to the
cohomology of $G$-modules, from now on write the group $A$ above as $M$,
require it to be commutative, and write its group operation additively. Such a
group extension determines a group homomorphism
\[ \alpha:G\to\Aut_{\cate{Grp}}(M),\qquad
	\alpha(g)(x)=exe^{-1},\qquad e\in\pi^{-1}(g)\ \text{chosen arbitrarily}. \]
Thus $M$ becomes a $G$-module through $\alpha$; this is the most basic
invariant of a group extension.

% segment-id: o014.li2.chapter6.low-degree-cohomology.p010
Given a $G$-module $M$, we wish to classify all extensions of $G$ by $M$
whose associated $\alpha$ is precisely the given $G$-module structure on $M$,
up to equivalence (or isomorphism) of extensions. An equivalence from an
extension $E$ to an extension $E'$ is defined in the evident way as a
commutative diagram of groups
\begin{equation}\label{eqn:group-ext-diagram}\begin{tikzcd}
	0 \arrow[r] & M \arrow[r] \arrow[equal,d] &
	E \arrow[r] \arrow[d,"\varphi"] & G \arrow[r] \arrow[equal,d] & 1 \\
	0 \arrow[r] & M \arrow[r] & E' \arrow[r] & G \arrow[r] & 1.
\end{tikzcd}\end{equation}
The reader is invited to verify that $\varphi$ in this diagram is
automatically a group isomorphism, and that an extension is split if and only
if it is equivalent to the semidirect product $M\rtimes G$.

% segment-id: o014.li2.chapter6.low-degree-cohomology.d003
\begin{definition}
	\index[sym1]{ExtGM@$\cate{Ext}(G,M)$}
	Define the category $\cate{Ext}(G,M)$ of all extensions of $G$ by the
	given $G$-module $M$; its morphisms are the commutative diagrams
	\eqref{eqn:group-ext-diagram}.
\end{definition}

% segment-id: o014.li2.chapter6.low-degree-cohomology.p011
This category is a groupoid: all its morphisms are invertible. Classifying
extensions amounts to describing $\cate{Ext}(G,M)$ up to equivalence of
categories. This divides into two subproblems:
\begin{itemize}
	\item describe $\Obj(\cate{Ext}(G,M))/\simeq$, that is, the equivalence
	classes of extensions, and identify the split extensions among them;
	\item describe the morphisms in $\cate{Ext}(G,M)$, that is, the
	isomorphisms between extensions.
\end{itemize}
We begin with the simpler second subproblem. First, an extension $E$ has a
simple class of automorphisms of the form $\Ad_m:e\mapsto mem^{-1}$, where
$m\in M$ and the group operation in $E$ is written multiplicatively; these
are called inner automorphisms. For abstract reasons, the inner automorphisms
form a normal subgroup of $\Aut(E)$, and the corresponding quotient group is
called the outer automorphism group of $E$.

% segment-id: o014.li2.chapter6.low-degree-cohomology.p012
To classify group extensions, we need to compute cohomology using the
normalized standard complex $\overline{C}(G,M)$ described in
Remark~\ref{rem:nor-cochain}. Accordingly, define
\begin{gather*}
	\overline{C}^n(G,M)\supset\overline{Z}^n(G,M)
	\supset\overline{B}^n(G,M),\\
	\Hm^n(G,M)\simeq\overline{Z}^n(G,M)/\overline{B}^n(G,M).
\end{gather*}
We have $\overline{C}^0(G,M)=C^0(G,M)=M$. The preceding discussion of
crossed homomorphisms has also shown that
$\overline{Z}^1(G,M)=Z^1(G,M)$.

% segment-id: o014.li2.chapter6.low-degree-cohomology.l001
\begin{lemma}\label{prop:factor-set-prep}
	Let $0\to M\to E\xrightarrow{\pi}G\to1$ be an object of
	$\cate{Ext}(G,M)$. Its automorphism group is canonically isomorphic to
	$\overline{Z}^1(G,M)$, while its outer automorphism group is canonically
	isomorphic to $\Hm^1(G,M)$. More explicitly, an element
	$z\in\overline{Z}^1(G,M)$ corresponds to the automorphism sending
	$e\in E$ to $z(\pi(e))e$.
\end{lemma}
\begin{proof}
	For the moment, write the group operation on $M$ multiplicatively. Take
	$E'=E$ in the commutative diagram \eqref{eqn:group-ext-diagram}. For every
	$e\in E$, there is a unique $z(e)\in M$ such that $\varphi(e)=z(e)e$.
	If $m\in M$, then
	$\varphi(me)=m\varphi(e)=mz(e)e=z(e)me$, so $z(e)$ depends only on the
	image of $e$ in $G$; henceforth we write it as a map $z:G\to M$.
	Conversely, every map $z:G\to M$ determines a map $\varphi:E\to E$ by
	$\varphi(e)=z(\pi(e))e$, and \eqref{eqn:group-ext-diagram} commutes at
	the level of sets.

	Consider a map $z:G\to M$ and the corresponding $\varphi$. Put
	$g_i=\pi(e_i)$ for $i=1,2$. Then
	\begin{align*}
		\varphi(e_1e_2)&=z(g_1g_2)e_1e_2,\\
		\varphi(e_1)\varphi(e_2)&=z(g_1)e_1z(g_2)e_2\\
		&=z(g_1)\underbracket{g_1\cdot z(g_2)}_{
		\text{the $G$-module structure on $M$}}e_1e_2.
	\end{align*}
	Thus $\varphi$ is a group homomorphism if and only if $z$ is a crossed
	homomorphism. It is clear that composition of group homomorphisms
	corresponds to addition of crossed homomorphisms. Now take $m\in M$.
	The corresponding inner automorphism is
	$e\mapsto mem^{-1}=mem^{-1}e^{-1}e$. Clearly
	$mem^{-1}e^{-1}\in M$, and in additive notation this element is
	$m-\pi(e)m$. This proves the assertion.
\end{proof}

% segment-id: o014.li2.chapter6.low-degree-cohomology.d004
\begin{definition}
	The groupoid $\tau^{\leq2}\overline{\cate{C}}(G,M)$ is defined as follows:
	its set of objects is $\overline{Z}^2(G,M)$, while its morphism sets are
	\[ \Hom(f,f'):=\{w\in\overline{C}^1(G,M):f'=d^1w+f\}. \]
	Composition of morphisms is defined by addition in
	$\overline{C}^1(G,M)$. Thus the isomorphism classes of objects correspond
	bijectively to the elements of $\Hm^2(G,M)$.
\end{definition}

% segment-id: o014.li2.chapter6.low-degree-cohomology.p013
The meaning of the notation is clear:
$\tau^{\leq2}\overline{\cate{C}}(G,M)$ is the category determined by the
truncated complex $\tau^{\leq2}\overline{C}(G,M)$.

% segment-id: o014.li2.chapter6.low-degree-cohomology.t003
\begin{theorem}\label{prop:factor-set}
	Let $M$ be a $G$-module. The category $\cate{Ext}(G,M)$ is equivalent to
	$\tau^{\leq2}\overline{\cate{C}}(G,M)$. In particular, equivalence
	classes of extensions of $G$ by $M$ correspond bijectively to elements
	of $\Hm^2(G,M)$, and the equivalence class of split extensions
	corresponds to $0\in\Hm^2(G,M)$.
\end{theorem}
\begin{proof}
	We first construct a functor
	$\tau^{\leq2}\overline{\cate{C}}(G,M)\to\cate{Ext}(G,M)$. Given a map
	$f:G^2\to M$, define a binary operation on the set $M\times G$ by
	\[ (m_1,g_1)\cdot(m_2,g_2):=
	\left(m_1+g_1\cdot m_2+f(g_1,g_2),g_1g_2\right). \]
	Associativity is equivalent to the identity
	\begin{equation}\label{eqn:grp-2-cocycle}
		f(g_1,g_2)+f(g_1g_2,g_3)
		=g_1\cdot f(g_2,g_3)+f(g_1,g_2g_3),
	\end{equation}
	that is, to $f\in Z^2(G,M)$. Furthermore, $(0,1_G)$ is the identity if
	and only if
	\[ f(g,1_G)=0=f(1_G,g), \]
	that is, if and only if $f\in\overline{Z}^2(G,M)$. Taking
	$(g_1,g_2,g_3)=(g,g^{-1},g)$ in \eqref{eqn:grp-2-cocycle} and using the
	last equality gives
	\begin{equation}\label{eqn:grp-2-cocycle-inv}
		f(g,g^{-1})=g\cdot f(g^{-1},g).
	\end{equation}

	Thus $M\times G$ acquires a monoid structure from
	$f\in\overline{Z}^2(G,M)$; denote this monoid by $E$. It is automatically
	a group, since \eqref{eqn:grp-2-cocycle-inv} gives
	\[ (0,g)^{-1}=\left(-f(g^{-1},g),g^{-1}\right),\qquad
	(m,1_G)^{-1}=(-m,1_G). \]
	It is easy to prove that the maps $m\mapsto(m,1_G)$ and $(m,g)\mapsto g$
	give a group extension $0\to M\to E\to G\to1$ that induces precisely the
	original $G$-action on $M$. Verify that $f=0$ corresponds to the
	semidirect product $E=M\rtimes G$.

	If, moreover, $w\in\overline{C}^1(G,M)$ and $f':=d^1w+f$, denote the
	corresponding extension by $E'$. The map
	\[ (m,g)\mapsto(m-w(g),g) \]
	gives an isomorphism from $E$ to $E'$. Indeed, the condition that this
	map preserve multiplication is equivalent to the identity
	\[ -w(g_1)-g_1\cdot w(g_2)+(d^1w)(g_1,g_2)=-w(g_1g_2), \]
	which is exactly the definition of $d^1w$.

	We have therefore obtained a functor
	$\tau^{\leq2}\overline{\cate{C}}(G,M)\to\cate{Ext}(G,M)$, and both
	categories are groupoids. Comparing automorphism groups of objects by
	Lemma~\ref{prop:factor-set-prep} shows that this functor is fully
	faithful. It remains to show that it is essentially surjective.

	Let $E$ be an object of $\cate{Ext}(G,M)$. Choose an arbitrary section
	$s:G\to E$ of the map $\pi:E\to G$ in the extension data; that is, a map
	satisfying $\pi s=\identity_G$. This is equivalent to choosing a
	representative of every $M$-coset in $E$. A section $s$ is called
	normalized if $s(1_G)=1_E$. For all $g_1,g_2\in G$, the elements
	$s(g_1g_2)$ and $s(g_1)s(g_2)$ have the same image under $\pi$. Hence
	there is a unique $f(g_1,g_2)\in M$ such that
	\[ s(g_1)s(g_2)=f(g_1,g_2)s(g_1g_2). \]

	It is easy to prove that associativity,
	$(s(g_1)s(g_2))s(g_3)=s(g_1)(s(g_2)s(g_3))$, translates into
	\eqref{eqn:grp-2-cocycle}, namely $f\in Z^2(G,M)$. If $s$ is normalized,
	then $f(1_G,g)=0=f(g,1_G)$ for every $g$, which is equivalent to
	$f\in\overline{Z}^2(G,M)$.

	For the chosen normalized section $s$, give $M\times G$ the group
	structure corresponding to $f$. Consider the bijection
	\[ M\times G\xrightarrow{1:1}E,\qquad(m,g)\mapsto ms(g). \]
	The preceding discussion shows immediately that this bijection preserves
	group multiplication and gives an isomorphism in $\cate{Ext}(G,M)$.
	This completes the proof of the equivalence.
\end{proof}

% segment-id: o014.li2.chapter6.low-degree-cohomology.p014
In the final part of the essential-surjectivity argument, changing the section
$s$ is equivalent to choosing an arbitrary map $h:G\to M$ and taking
$s'(g):=h(g)s(g)$. The corresponding map $f':G^2\to M$ becomes
\begin{align*}
	f'(g_1,g_2)&=\left(h(g_1)+g_1\cdot h(g_2)-h(g_1g_2)\right)+f(g_1,g_2)\\
	&=(d^1h)(g_1,g_2)+f(g_1,g_2).
\end{align*}
If one also requires $h(1_G)=0$, then
$d^1h\in\overline{B}^2(G,M)$. Historically, this observation was the
starting point for the connection between extensions and $\Hm^2(G,M)$; it
also supplies a group-extension interpretation for the definitions of
$\overline{Z}^2(G,M)$ and $\overline{B}^2(G,M)$, although this was left
implicit in the proof.

% segment-id: o014.li2.chapter6.low-degree-cohomology.r001
\begin{remark}\label{rem:Ext-incarnation}
	Theorem~\ref{prop:factor-set} says that the complex
	$\tau^{\leq2}\overline{C}(G,M)$ is a linear incarnation of the category
	of group extensions $\cate{Ext}(G,M)$: the term in degree $2$ describes
	objects, the term in degree $1$ describes morphisms, and the term in
	degree $0$ remains to be interpreted. For
	$m\in M=\overline{C}^0(G,M)$, there is an inner automorphism of the group
	extension $\Ad_m:e\mapsto mem^{-1}$, which is trivial if and only if
	$m\in M^G=\Hm^0(G,M)$. From the viewpoint of inner automorphisms, we may
	regard the elements of $M$ as 2-morphisms in $\cate{Ext}(G,M)$: for
	$m\in M$ and
	$\varphi_1,\varphi_2\in\Hom_{\cate{Ext}(G,M)}(E,E')$,
	\[ \text{interpret the 2-cell diagram}\;
	\begin{tikzcd}[column sep=large]
		E \arrow[bend left=60,r,"\varphi_1",""' {name=U}]
		\arrow[bend right=60,r,"\varphi_2"',"" {name=D}] &
		E' \arrow[Rightarrow,from=U,to=D,"m"]
	\end{tikzcd}
	\;\text{as}\;\varphi_2=\Ad_m\varphi_1. \]

	Notice that $\Ad_m\varphi_1=\varphi_1\Ad_m$ and
	$\Ad_{m_1}\Ad_{m_2}=\Ad_{m_1+m_2}$. Thus $\Hm^0(G,M)$ becomes the
	“2-automorphism group” of every morphism $\varphi$. In this sense, the
	complex $\tau^{\leq2}\overline{C}(G,M)$ may be called a linear
	incarnation of the 2-category $\cate{Ext}(G,M)$. See
	\cite[\S 3.5]{Li1} for the definition of a 2-category.
\end{remark}

% segment-id: o014.li2.chapter6.low-degree-cohomology.p015
Other aspects of group extensions will be discussed in
\S\ref{sec:group-ext-revisited}. Before that, we need to introduce more
operations on group homology and cohomology.
